%! AP Statistics — Unit 4, Lesson 4.5 — Warm-Up
\documentclass[10pt]{article}
\usepackage{apstats-article}
\usepackage{apstats-boxes}

\begin{document}

\pageheader{Unit 4, Lesson 4.5}{Warm-Up}
\namedateperiod

\vspace{0.3in}

% ── SECTION 1: Spiral Review ──────────────────────────────────────────────────
{\large\bfseries\color{navy} Part 1 \quad Spiral Review}

\vspace{0.12in}

\begin{enumerate}

    \item Use the two-way table below, which shows data for 120 athletes.

    \vspace{6pt}
    \renewcommand{\arraystretch}{1.4}
    \begin{tabularx}{\linewidth}{@{} l X X r @{}}
     & \textbf{Trains Daily} & \textbf{Trains Weekly} & \textbf{Total} \\
    \hline
    \textbf{Runs}     & 24 & 16 & 40 \\
    \textbf{Swims}    & 18 & 22 & 40 \\
    \textbf{Cycles}   & 20 & 20 & 40 \\
    \hline
    \textbf{Total}    & 62 & 58 & 120 \\
    \end{tabularx}
    \vspace{6pt}
    One athlete is selected at random.
    \begin{enumerate}[label=(\alph*), itemsep=6pt]
        \item What is $P(\text{Runs})$? \blank{5cm}
        \item What is $P(\text{Trains Daily})$? \blank{5cm}
        \item What is $P(\text{Runs AND Trains Daily})$, i.e.\ $P(\text{Runs} \cap \text{Trains Daily})$?\\
              \blank{5cm}
        \item Are ``Runs'' and ``Swims'' mutually exclusive? Justify briefly.\\
              \writeline
    \end{enumerate}

    \vspace{0.18in}

    \item Recall from Lesson 4.4: two events $A$ and $B$ are mutually exclusive if
    $P(A \cap B) = $ \blank{2cm}.

    \vspace{0.14in}

    \item A card is drawn from a standard 52-card deck.
    \begin{enumerate}[label=(\alph*), itemsep=6pt]
        \item $P(\text{King}) = $ \blank{3cm}
        \item $P(\text{King} \cap \text{Heart}) = $ \blank{3cm}
        \item If you already know the card is a Heart, how many cards are left
              to consider? \blank{2cm} \quad How many of those are Kings? \blank{2cm}
    \end{enumerate}

\end{enumerate}

\vspace{0.3in}

% ── SECTION 2: Looking Ahead ──────────────────────────────────────────────────
{\large\bfseries\color{navy} Part 2 \quad Looking Ahead}

\vspace{0.12in}

\noindent Today we ask: \textit{``If we already know one event occurred, how does that
change the probability of another?''}

\vspace{0.14in}

\begin{tcolorbox}[colback=greenbg, colframe=greenacc, boxrule=0.8pt, arc=2mm,
    left=4mm, right=4mm, top=3mm, bottom=3mm]
    \small
    \begin{itemize}[leftmargin=1.3em, itemsep=8pt]
      \item Conditional probability restricts the sample space to the ``given'' event.
      \item The formula is $P(A|B) = P(A \cap B) / P(B)$.
      \item From a table: numerator = joint cell; denominator = conditioning row/column total.
    \end{itemize}
\end{tcolorbox}

\vspace{0.2in}

\noindent\textcolor{slate}{\small Write one question about today's lesson or
something you want to understand better:}\\[8pt]
\writeline\\[2pt]\writeline

\end{document}
